%----------------------------------
\section{Data Science Role}
\label{sec:Data Science Role}
%----------------------------------

%%----------------------------------
\subsection{Job Description and Role Identification}
\label{subsec:Job Description and Role Identification}
%----------------------------------

The role selected is \textbf{Machine Learning Engineer Graduate (Trust and Safety
Engineering) -- 2027 Start} at \textbf{TikTok}, based in Sydney, New South Wales, and
open to international applicants. Applications close 30 August 2026. The full
advertisement is reproduced in Appendix~\ref{sec:Job Listing}.

The advertisement places the engineer inside TikTok's Trust \& Safety team, which
describes itself as having ``invested heavily in human and machine-based moderation to
remove harmful content quickly and often before it reaches our general community''. The
listed responsibilities are explicitly build-oriented: to ``develop and build up
highly-scalable classifiers, tools, models and algorithms leveraging cutting-edge machine
learning, computer vision and data mining technologies'', to ``improve our trust and
safety strategy and work on model iterations'', and to ``develop and implement innovative
machine learning algorithms to manage business risks''. Minimum qualifications are a
Master's or PhD in computer science or a related engineering field, solid knowledge of at
least one of machine learning, pattern recognition, natural language processing, data
mining or computer vision, and a firm understanding of data structures and algorithms.

The professional problem behind the role is one of \emph{scale under a fixed review
budget}. A platform of TikTok's size receives far more user-generated content each day
than any moderation team can read, yet removing content automatically and without recourse
is unacceptable to creators and regulators alike. Neither exhaustive human review nor
fully autonomous removal is viable. The engineering problem is therefore triage: rank
incoming items by the probability that they violate policy, so that a finite number of
human moderators spend their attention where it changes an outcome. This is the problem
the remainder of this report addresses, and it is what makes the role a machine learning
\emph{engineering} role rather than an analytics one --- the deliverable is a ranked queue
running in production, monitored and re-trained, not a slide deck.

\textbf{Job listing URL:} \url{https://au.gradconnection.com/employers/tiktok/jobs/tiktok-machine-learning-engineer-graduate-trust-and-safety-engineering-2027-start-5/}

%%----------------------------------
\subsection{Industry Context}
\label{subsec:Industry Context}
%----------------------------------

Short-form video platforms cannot human-review millions of daily uploads. The trust and
safety field uses data science to detect hate speech, harassment, spam and coordinated
abuse, to rank what scarce moderators see first, and to monitor whether enforcement stays
consistent across languages, dialects and regions.

\wordcount{45}

%%----------------------------------
\subsection{Role Focus}
\label{subsec:Role Focus}
%----------------------------------

Prediction, delivered as production software. The advertisement asks the engineer to
``develop and build up highly-scalable classifiers, tools, models and algorithms'' and to
``work on model iterations'' to ``manage business risks''. Insight generation is
secondary; the deliverable is a deployed classifier that ranks content for moderator
review.

\wordcount{47}

%%----------------------------------
\subsection{Value Creation}
\label{subsec:Value Creation}
%----------------------------------

Automated triage lets a fixed moderation team cover an expanding content stream, so
harmful material is removed before it spreads. Better ranking also protects creators from
wrongful removal and advertisers from unsafe placement, while producing the enforcement
evidence regulators such as Australia's eSafety Commissioner require~\cite{esafety2022bosa}.

\wordcount{46}

%%----------------------------------
\subsection{Values and Culture Alignment}
\label{subsec:Values and Culture Alignment}
%----------------------------------

The team frames its mission as keeping people ``safe and empowered to create'', which
matches how I want to use machine learning: reducing harm rather than maximising
engagement. I value measured claims over confident ones, and this role's emphasis on
experiments, validation and iteration rewards that habit.

\wordcount{47}

%%----------------------------------
\subsection{Diversity and Inclusion}
\label{subsec:Diversity and Inclusion}
%----------------------------------

Moderation models inherit their annotators' judgements. The Wikipedia corpus used here
ships annotator demographics precisely because who labels the data shapes what counts as
toxic~\cite{wulczyn2017ex}. Diverse engineering teams are likelier to disaggregate error
rates by dialect and community, catching the over-flagging that silences minority
speakers~\cite{sap2019risk}.

\wordcount{47}

%%----------------------------------
\subsection{Cover Letter}
\label{subsec:Cover Letter}
%----------------------------------

A cover letter tailored to this advertisement is presented as a stand-alone letter in
Appendix~\ref{sec:Cover Letter}, formatted so that it could be sent unaltered with a job
application.

\wordcount{348 (see Appendix~\ref{sec:Cover Letter})}
